\section*{Chapitre \ref{ch:ll:multi-socket-machines-and-interconnects} --- Machines multi-sockets et interconnexions}

\begin{solution}{exo:ll:multi-socket-machines-and-interconnects:1}
$16 \times 8 \times 128/130/8 = 15.75$~Go/s dans chaque sens~; deux ports de 25~Gb/s demandent 6.25~Go/s~: oui, avec de la marge.
\end{solution}

\begin{solution}{exo:ll:multi-socket-machines-and-interconnects:2}
$\lceil 300/64 \rceil = 5$ lignes si le paquet commence en début de ligne, 6 sinon.
\end{solution}

\begin{solution}{exo:ll:multi-socket-machines-and-interconnects:3}
La carte est sur le nœud~1 et la stratégie sur le nœud~0~: accès distant. $(3 + 1)(138 - 59) = \qty{316}{\nano\second}$ par message.
\end{solution}

\begin{solution}{exo:ll:multi-socket-machines-and-interconnects:4}
L'\omterm{def:ll:multi-socket-machines-and-interconnects:first}{allocation au premier accès} (first touch) a placé les pages de l'anneau sur le nœud~0, donc chaque écriture depuis le nœud~1 traverse
l'interconnexion pour atteindre la mémoire ou les caches du nœud~0, et le lecteur paie à nouveau s'il se trouve sur le
nœud~1. Allouer et mettre à zéro l'anneau dans le thread de flux après l'avoir épinglé, ou lier sa mémoire au nœud~1.
\end{solution}

\begin{solution}{exo:ll:multi-socket-machines-and-interconnects:5}
Deux CPU virtuels voisins, par exemple 0 et 1~: un passage de témoin d'environ 40~nanosecondes dans un sens sur la
matrice engagée, contre une médiane d'environ 190 pour les autres paires, soit quatre à cinq fois plus. Ce choix ne tient
que tant que la correspondance de l'hyperviseur reste la même.
\end{solution}

\begin{solution}{exo:ll:multi-socket-machines-and-interconnects:6}
Le jumeau partage les unités d'exécution du cœur, les caches de premier niveau et les prédicteurs, donc tout travail sur
celui-ci ralentit le thread chaud et augmente sa \omterm{def:ll:where-latency-comes-from:tail}{gigue} (jitter). Sur le montage, les trois threads chauds laissent les CPU 29, 30 et 31 inactifs~: trois des seize CPU du nœud~1.
\end{solution}

\begin{solution}{exo:ll:multi-socket-machines-and-interconnects:7}
Il faut regrouper sur un même cœur les deux threads qui échangent le moins de données par message et se disputent le
cœur le moins souvent~: la passerelle et le constructeur de carnet sont de mauvais candidats, puisqu'ils interviennent tous deux à chaque message~; la
paire qui n'échange de données que par l'intermédiaire d'un autre thread, par exemple le constructeur de carnet et la passerelle lorsque la stratégie se trouve entre eux, est le moins mauvais choix.
\texttt{check} ne doit accepter qu'une paire de rôles déclarée sur un même cœur physique, et rien d'autre.
\end{solution}

\begin{solution}{exo:ll:multi-socket-machines-and-interconnects:8}
\texttt{taskset} épingle le thread mais ne dit rien de la carte~: personne n'a vérifié que le CPU~5 se trouve sur le nœud
de la carte, le choix casse lorsque le serveur ou l'emplacement change, la mémoire du processus a pu être touchée pour la
première fois ailleurs, et les threads créés ensuite héritent du masque au lieu d'obtenir leurs propres cœurs. Il faut dériver le
placement de la topologie et le vérifier au démarrage.
\end{solution}

\begin{solution}{pb:ll:multi-socket-machines-and-interconnects:1}
\begin{enumerate}
\item \path{/sys/bus/pci/devices/<address>/numa_node} pour la carte (trouvée à partir de
\path{/sys/class/net/<name>/device}), et \path{/sys/devices/system/node/node*/cpulist} pour les CPU.
\item CPU 16 à 31~; les jumeaux du CPU~17 sont 17 et 25.
\item Flux sur 21, stratégie sur 22, passerelle sur 23~; CPU 29, 30 et 31 inoccupés.
\item Le noyau, les minuteries et de nombreux pilotes privilégient le CPU~0~; le réserver comme cœur de service (housekeeping)
éloigne ce travail du chemin critique (hot path).
\item $(2 + 2)(138 - 59) = \qty{316}{\nano\second}$.
\item $316/2\,500 = 12.6~\%$.
\item Chaque accès au carnet tombe désormais sur la mémoire du nœud~0 à travers l'interconnexion~: les threads sont
locaux à la carte, mais leurs données sont distantes.
\item Le carnet~: il effectue de nombreux accès par message, dépendants les uns des autres (parcours d'arbre ou d'échelle), alors que les lignes de la carte sont peu nombreuses.
\item Une ligne dans chaque sens~: $2 \times 59 = \qty{118}{\nano\second}$ au sein d'un socket, $2 \times 138 = \qty{276}{\nano\second}$
entre sockets.
\item Environ 40~nanosecondes dans un sens entre CPU voisins et environ 190 entre les autres, d'après la
mesure engagée.
\item Il faut choisir la paire au passage de témoin mesuré le plus faible, parmi les CPU du nœud de la carte, et la vérifier après chaque démarrage.
\item L'hyperviseur répartit les CPU virtuels sur les cœurs physiques comme il l'entend, et cette correspondance peut
changer d'un démarrage à l'autre ou pendant l'exécution.
\item \textbf{Résultat nommé.} Le mauvais socket ajoute environ \qty{316}{\nano\second} par message, soit 12.6~\% du
budget fil à fil de \qty{2.5}{\micro\second}, avant de compter les données dont la première écriture a eu lieu sur le mauvais nœud.
\item Lire la topologie~; calculer le plan~; épingler chaque thread chaud avant qu'il n'alloue~; allouer et toucher sa
mémoire depuis ce thread~; éloigner les interruptions des autres périphériques des cœurs chauds~; vérifier le plan au démarrage et refuser de trader s'il échoue.
\item Relire le CPU de chaque thread (\path{/proc/<pid>/task/*/stat} ou \texttt{sched\_getcpu} journalisé par le thread) et
le nœud de chaque région (\path{/proc/<pid>/numa_maps}), puis comparer le tout au plan.
\item Le plan est recalculé à partir de la topologie et le processus suit la carte~; une vérification au démarrage le
détecte si personne ne le fait.
\item Ceux des disques, des autres cartes réseau et des minuteries~; les interruptions propres à la carte, si elle en
utilise, vont aux cœurs de service du nœud.
\item Le micrologiciel n'a pas signalé le nœud de la carte~: le plan retombe sur une valeur par défaut et doit le
signaler~; un humain doit confirmer ce choix.
\item Le plan de placement (la topologie, le CPU et le nœud de chaque thread), avec les versions de la machine et des logiciels.
\item Parce que l'endroit où se trouvent un thread et ses données détermine le nombre de traversées de l'interconnexion que chaque message subit.
\end{enumerate}
\end{solution}

\begin{solution}{iq:ll:multi-socket-machines-and-interconnects:1}
Chaque socket possède sa propre mémoire~; la mémoire rattachée à un autre socket est plus lente à atteindre à travers
l'interconnexion. Un processus sensible à la latence doit garder ses threads chauds, leur mémoire et la carte réseau sur un même nœud.
\iqlookfor{mémoire locale contre mémoire distante, et conséquence pour le placement.}
\end{solution}

\begin{solution}{iq:ll:multi-socket-machines-and-interconnects:2}
Linux place une page sur le nœud du CPU qui l'écrit en premier~; si le thread principal initialise les données au
démarrage et que le worker tourne sur un autre nœud, les données du worker sont distantes. Il faut épingler d'abord, puis toucher, ou lier explicitement.
\iqlookfor{que c'est l'initialisation, et non l'allocation, qui décide du placement.}
\end{solution}

\begin{solution}{iq:ll:multi-socket-machines-and-interconnects:3} \sloppy
Trouver son adresse PCI à partir de \path{/sys/class/net/<if>/device}, puis lire \texttt{numa\_node} et \texttt{local\_cpulist}
sous \path{/sys/bus/pci/devices/<address>/}~; \texttt{lstopo} la montre graphiquement.
\iqlookfor{sysfs, et le fait que $-1$ signifie inconnu.}
\end{solution}

\begin{solution}{iq:ll:multi-socket-machines-and-interconnects:4}
Sur le même socket, idéalement sur des cœurs qui partagent un niveau de cache, chacun sur son propre cœur physique, les jumeaux étant inoccupés, sur le
nœud de la carte réseau~: un passage de témoin entre sockets coûte plus de deux fois plus.
\iqlookfor{latence de cœur à cœur comme coût d'un passage de témoin.}
\end{solution}

\begin{solution}{iq:ll:multi-socket-machines-and-interconnects:5}
La carte écrit les paquets dans la mémoire de l'hôte d'elle-même, sans que le processeur les recopie~; avec l'\omterm{def:ll:multi-socket-machines-and-interconnects:dca}{accès direct au cache}, ils
atterrissent dans le cache de dernier niveau du socket de la carte, de sorte que le cœur qui les lit trouve ses données en cache plutôt qu'en mémoire, à condition que
le cœur se trouve sur ce socket.
\iqlookfor{transferts pilotés par le périphérique et leur localité.}
\end{solution}

\begin{solution}{iq:ll:multi-socket-machines-and-interconnects:6}
Au démarrage, lire la topologie~; dériver le CPU de chaque thread à partir du nœud de la carte par une règle fixe (un cœur physique chacun,
jumeaux inoccupés, CPU 0 exclu)~; épingler chaque thread avant qu'il n'alloue~; allouer et toucher sa mémoire depuis lui ou la lier~;
vérifier le placement à partir de \path{/proc} et refuser de trader en cas d'écart~; journaliser le plan avec chaque mesure.
\iqlookfor{placement dérivé plutôt que saisi à la main, et une vérification qui refuse de démarrer.}
\end{solution}
